\subsection{Descent of groupoids}

Let X and Y be topological spaces, and let \(f\colon\mathrm{X}\to\mathrm{Y}\) be a continuous map. Let \(p_{1}\) and \(p_{2}\) be the two projections of \(\mathrm{X}\times_{\mathrm{Y}}\mathrm{X}\) into X, \(\operatorname{Coeg}(f)\) the coequalizer groupoid of the pair \((\varpi(p_{1}),\varpi(p_{2}))\) of groupoid morphisms from \(\varpi(\mathrm{X}\times_{\mathrm{Y}}\mathrm{X})\) into \(\varpi(\mathrm{X})\), and \(\gamma\colon\varpi(\mathrm{X})\to\operatorname{Coeg}(f)\) the canonical groupoid morphism (II, p.~199, def. 2). Since \(f\circ p_{1}=f\circ p_{2}\) , we have \(\varpi(f)\circ\varpi(p_{1})=\varpi(f)\circ\varpi(p_{2})\) . It then follows from the universal property of coequalizers (II, p.~199, prop. 3) that there exists a unique morphism of groupoids \(\varpi'(f)\colon\operatorname{Coeg}(f)\to\varpi(\mathrm{Y})\) such that \(\varpi(f)=\varpi'(f)\circ\gamma\) .

The set of vertices of Coeg(f) is the quotient set of the set X = Vert( \(\varpi(X)\) ) by the equivalence relation defined by f (II, p.~200, remark 1). It is identified with f(X).

PROPOSITION 8. --- Let X and Y be nonempty topological spaces and let f be a continuous map from X to Y. Suppose that X is locally path-connected, Y is connected, and f is strict and surjective. Then the groupoid Coeg(f) is transitive.

Let \(\Gamma\) be the quiver whose set of vertices is X and whose set of arrows is the disjoint union of \(\mathrm{Fl}(\varpi(\mathrm{X}))\) and \(\mathrm{X}\times_{\mathrm{Y}}\mathrm{X}\), the source and target maps being those of \(\varpi(\mathrm{X})\) on \(\mathrm{Fl}(\varpi(\mathrm{X}))\) and the maps \(p_1\) and \(p_2\) on \(\mathrm{X}\times_{\mathrm{Y}}\mathrm{X}\) . By the definition of the frame of the pair \((\varpi(\mathrm{pr}_1), \varpi(\mathrm{pr}_2))\) (II, p.~185, def. 3) and remark 2 of II, p.~200, the set of orbits of \(\operatorname{Coeg}(f)\) is identified with the set of connected components of the quiver \(\Gamma\) . Since the space X is not empty, it suffices to prove that the graph \(\Gamma\) is connected.

The connected components of \(\Gamma\) are saturated for the equivalence relation “there exists a path joining \(x\) to \(x'\) ”, hence are open in X, because X is locally path-connected. They are then also closed. They are also saturated for the equivalence relation R defined by f . By hypothesis, the map f induces a homeomorphism from X/R onto Y, so that the image under f of every connected component of \(\Gamma\) is an open and closed subset of Y, hence is equal to Y since the space Y is assumed to be connected.

Let C be a connected component of \(\Gamma\) , let \(x\) be a point of X. By the preceding, there exists a point \(x' \in \mathrm{C}\) such that \(f(x') = f(x)\) . By definition of the quiver \(\Gamma\) , we have \(x' \in \mathrm{C}\) . One thus has \(\mathrm{C} = \mathrm{X}\) and the quiver \(\Gamma\) is therefore connected.

PROPOSITION 9. --- Let X and Y be topological spaces and let f be a continuous map from X to Y. Suppose that the space X is locally path-connected, that the space Y is deloopable, and that the map f is strict and surjective. Then, the groupoid \(\varpi(\mathrm{Y})\) is generated by the image of \(\varpi(\mathrm{X})\) under \(\varpi(f)\) .

We may assume that the spaces X and Y are nonempty. Let G denote the subgroupoid of \(\varpi(\mathrm{Y})\) generated by the image of \(\varpi(\mathrm{X})\) under \(\varpi(f)\) . As the map f is surjective, the set of vertices of G is equal to Y. By prop. 8, \(\operatorname{Coeg}(f)\) is a transitive groupoid. Since \(\varpi'(f)\) induces the identity on the vertices, the image of \(\operatorname{Coeg}(f)\) under \(\varpi'(f)\) is a transitive subgroupoid of \(\varpi(\mathrm{Y})\) . The image of \(\varpi(\mathrm{X})\) under \(\gamma\) generates \(\operatorname{Coeg}(f)\) (II, p.~200, corollary); since we have \(\varpi(f) = \varpi'(f) \circ \gamma\) , the groupoid G is transitive.

Let \(y_{0}\) be a point of Y and let H denote the subgroup \(G_{y_{0}}\) of \(\pi_{1}(Y, y_{0})\) . By th. 1 of IV, p.~342, there exists a connected covering (T, p) of Y and a point \(t_{0}\) of the fiber \(T_{y_{0}}\) whose fixator is H.

If \(x\) is a point of X, the set \(\mathrm{Fl}_{y_0,f(x)}(\mathrm{G})\) is nonempty, since G is transitive. For \(u\in \mathrm{Fl}_{y_0,f(x)}(\mathrm{G})\) , the point \(t_0\cdot u\) is a point of the fiber \(\mathrm{T}_{f(x)}\) , independent of \(u\) since the group H, the isotropy group of G at \(y_0\) , fixes \(t_0\) . Denote by \(\sigma (x)\) this point; denote by \(\sigma \colon \mathrm{X}\to \mathrm{T}\) the map thus defined and \(s\colon \mathrm{X}\to \mathrm{X}\times_{\mathrm{Y}}\mathrm{T}\) the map \(x\mapsto (x,\sigma (x))\) . By construction, the map s is compatible with the canonical actions of \(\varpi (\mathrm{X})\) on X and \(\mathrm{X}\times_{\mathrm{Y}}\mathrm{T}\) . It then follows from Lemma 1 of III, p.~312, that s is continuous. The map \(\sigma\) is therefore continuous; moreover, it is compatible with the equivalence relation defined by f since \(\sigma (x)\) depends only on \(f(x)\) . Since f is strict and surjective, there exists a unique continuous map \(\overline{\sigma}\colon \mathrm{Y}\to \mathrm{T}\) such that \(\overline{\sigma}\circ f = \sigma\) , so that the covering T admits a section. Since T is connected, the map \(p\colon \mathrm{T}\to \mathrm{Y}\) is a homeomorphism (I, p.~31, cor. 4 of prop. 6) and one has \(\mathrm{H} = \pi_1(\mathrm{Y},y_0)\) , whence \(\mathrm{G}_{y_0} = \pi_1(\mathrm{Y},y_0)\) . Since G is transitive, it follows that \(\mathrm{G} = \varpi (\mathrm{Y})\) .

PROPOSITION 10. --- Let X and Y be topological spaces and let f be a continuous map from X to Y. Suppose that the space X is deloopable, that the space \(X\times_{Y}X\) is locally path-connected, and that every descent datum relative to f on a covering of X is effective and the quotient space is a covering of Y. Then, the groupoid morphism \(\varpi'(f)\) from Coeg(f) into \(\varpi(Y)\) is injective.

Under the hypotheses of the proposition, \(f\) is strict (IV, p.~394, remark); one may moreover assume that it is surjective.

Lemma 2. --- Keep the notations and hypotheses of the proposition. Let T be a set equipped with an action of the groupoid Coeg(f) relative to a map q: T → Y. Then there exists a unique topology on T for which q makes T a covering of Y such that for every path class with origin c in X and every point t in T, one has \(t \cdot \gamma(c) = t \cdot f_{*}(c)\) . In particular, the action of Coeg(f) on T factors through an action of the groupoid \(\varpi(Y)\) .

The uniqueness of such a topology on T follows from Lemma 1 (III, p.~312); let us prove its existence. Endow the set \(X\times_{Y}T\) with an action of \(\varpi(\mathrm{X})\) by setting \((x,t)\cdot u=(x\cdot u,t\cdot\gamma(u))\) for all \((x,t)\in\mathrm{X}\times_{\mathrm{Y}}\mathrm{T}\) and every arrow u with source x in \(\varpi(\mathrm{X})\). Since the space X is deloopable, this action is without local monodromy (III, p.~313, remark). There therefore exists on the set \(X\times_{Y}T\) a unique topology that makes the X-space \((\mathrm{X}\times_{\mathrm{Y}}\mathrm{T},\mathrm{pr}_{1})\) a covering of X such that the canonical action of \(\varpi(\mathrm{X})\) on this covering coincides with the given action (III, p.~313, prop. 3). Let us denote this covering by \((Z,p)\). Let us show that the map \(\tau: ((x_{1},t),x_{2})\mapsto(x_{2},t)\) from \(Z\times_{Y}X\) into Z is a descent datum on Z relative to the map f. It suffices to verify that it is continuous, the other conditions of Definition 1 of IV, p.~382, being evident.

The map \(\tau\) is a lifting of the map \(\mathrm{pr}_2\colon Z\times_Y X\to X\) to the covering \(Z\) of \(X\) . Since the space \(X\times_Y X\) is locally path-connected, the space \(Z\times_Y X\) , which is a covering of it, is also locally path-connected. By the corollary (III, p.~269), to show that the map \(\tau\) is continuous, it suffices to show that it is continuous by paths.

So let \(\widetilde{c} = ((c,g),c')\) be a path in \(Z\times_{Y}X\) and let us show that the map \(\tau \circ \widetilde{c}\colon t\mapsto (c'(t),g(t))\) from I to Z is continuous. For every \(s\in [0,1]\) , denote by \(c_s\) and \(c_s'\) the paths in X defined by \(t\mapsto c(st)\) and \(t\mapsto c'(st)\) . For all \(s\in [0,1]\) , we thus have \(c(s) = c(0)\cdot [c_s]\) , \(c'(s) = c'(0)\cdot [c_s']\) and \((c,g)(s) = (c(0),g(0))\cdot [c_s]\) , where \([u]\) denotes the strict homotopy class of a path u (III, p.~304, remark). The map \(t\mapsto (c_s(t),c_s'(t))\) is a path in \(X\times_{Y}X\) ; by definition of the coequalizer Coeg(f), we therefore have \(\gamma ([c_s]) = \gamma ([c_s'])\) and \(g(s) = g(0)\cdot \gamma ([c_s]) = g(0)\cdot \gamma ([c_s'])\) . By definition of the action of \(\varpi (X)\) on Z, we therefore have

\[
\begin{array}{r} (\tau \circ \widetilde {c}) (s) = (c ^ {\prime} (s), g (s)) = (c ^ {\prime} (0) \cdot [ c _ {s} ^ {\prime} ], g (0) \cdot \gamma ([ c _ {s} ^ {\prime} ])) \\ = (c ^ {\prime} (0), g (0)) \cdot [ c _ {s} ^ {\prime} ] = (\tau \circ \widetilde {c}) (0) \cdot [ c _ {s} ^ {\prime} ]. \end{array}
\]

This shows that \(\tau \circ \widetilde{c}\) is a continuous lift to Z of the path \(c'\) (loc. cit.). Consequently, the map \(\tau\) is continuous by paths, hence continuous; it is a descent datum relative to \(f\) on the X-space Z.

Denote by \(R_{\tau}\) the equivalence relation defined by the descent datum \(\tau\) . Since the map f is surjective, the map \(pr_{2}: Z \to T\) induces, by passing to the quotient, a bijection from \(Z/R_{\tau}\) onto T. Endow T with the topology deduced from that of \(Z/R_{\tau}\) by transport of structure, so that \((T,q)\) is a Y-space. By hypothesis, T is therefore a covering of Y and the diagram

\[
\begin{array}{c} \mathrm{Z} \xrightarrow {\mathrm{pr} _ {2}} \mathrm{T} \\ p \Big \downarrow \qquad \qquad \qquad \qquad \Big \downarrow q \\ \mathrm{X} \xrightarrow {f} \mathrm{Y} \end{array}
\]

is a Cartesian square.

Let \((x,t)\) be a point of Z, let c be a path with origin x in X, and let \(\widetilde{c}\) be the continuous lift of c to Z with origin \((x,t)\) . The path \(pr_{2} \circ \widetilde{c}\) is the lift with origin t of the path \(f \circ c\) of Y, so that \(t \cdot \gamma([c]) = t \cdot [f \circ c] = t \cdot f_{*}([c])\) . This concludes the proof of the lemma.

Let us now prove the proposition. Let \(u\) and \(v\) be two arrows of Coeg( \(f\) ) whose images in \(\varpi(Y)\) are equal. The groupoid morphism \(\varpi'(f)\) being the identity on the vertex sets, the arrows \(u\) and \(v\) have the same origin and the same target. Let us denote by \(y\) the origin of \(u\) ; let T be the set of arrows of Coeg( \(f\) ) with source y and let \(q: \mathrm{T} \to \mathrm{Y}\) the restriction to T of the target map. The groupoid Coeg( \(f\) ) acts by right composition on the set T, relative to \(q\) . By Lemma 2, the actions of \(u\) and \(v\) on T are identical. We therefore have \(u = e_y \cdot u = e_y \cdot v = v\) . This proves that the groupoid morphism \(\varpi'(f)\) is injective.

THEOREM 1. --- Let X and Y be topological spaces, let \(f: X \to Y\) be a continuous and surjective map. Suppose the spaces X and Y are deloopable. Finally, suppose that one of the following properties is satisfied:

\begin{enumerate}
\def\labelenumi{(\roman{enumi})}
\item
  The map f is proper, separated, with locally connected fibers, and the space \(X\times_{Y}X\) is locally path-connected;
\item
  The map f is proper, separated, with finite fibers, the diagonal \(\Delta_{X}\) is open in \(X \times_{Y} X\) and its complement is locally connected;
\item
  The map f is open and has the path lifting property.
\end{enumerate}

Then, the morphism of groupoids \(\varpi'(f)\) is an isomorphism from the groupoid Coeg(f) onto the Poincaré groupoid \(\varpi(Y)\) .

Let us first note that under these hypotheses, the map \(f\) is surjective and strict (I, p.~18, example 2). Moreover, every descent datum relative to \(f\) on a covering of X is effective, and the quotient space is a covering of Y; this follows from IV, p.~393, corollary 2 of prop. 4 under hypotheses (i) and (ii), and from prop. 7 of IV, p.~394 under hypothesis (iii). By prop. 10, the morphism of groupoids \(\varpi'(f)\) is therefore injective, and its image is a subgroupoid of \(\varpi(Y)\) .

By prop. 9 of IV, p.~395, this image is equal to \(\varpi(\mathrm{Y})\) under hypotheses (i) and (ii), but also under hypothesis (iii) since the groupoid morphism \(\varpi(f)\) is then surjective.

Therefore, \(\varpi'(f)\) is an isomorphism.

Examples. --- Here are two examples where the hypotheses of Theorem 1 are satisfied.

\begin{enumerate}
\def\labelenumi{\arabic{enumi})}
\item
  Let Y be a deloopable topological space. Let \((\mathrm{A}_{i})_{i\in\mathrm{I}}\) be a locally finite cover of Y by closed sets. Suppose that, for every \(i\in I\) , the space \(A_{i}\) is deloopable and that, for every pair \((i,j)\in\mathrm{I}\times\mathrm{I}\) , the space \(A_{i}\cap A_{j}\) is locally path-connected. One may take for X the sum space of the family \((\mathrm{A}_{i})_{i\in\mathrm{I}}\) and for \(f:X\to Y\) the map deduced from the family of canonical injections.
\item
  Let G be a discrete group acting properly on a deloopable topological space X; set Y = X/G and let \(f: X \to Y\) be the canonical map. It is open (TG, III, p.~10, lemma 2) and has the path-lifting property by virtue of Theorem 4 of III, p.~287. The space Y is deloopable according to IV, p.~349, prop. 8, b). By hypothesis, the map from G × X to \(X\times X\) given by \((g, x) \mapsto (g \cdot x, x)\) is proper, hence strict (I, p.~18, example 2), and its image is \(X\times_{Y}X\). It then follows from III, p.~261, prop. 8 that \(X\times_{Y}X\) is locally path-connected.
\end{enumerate}

